\section{Contributions, plan and risks}\label{sec:plan}
\textbf{Expected contributions.} (1) A quantified answer to how sensitive FRD, MMD and JSD are to ray drop, compared with real-data mIoU. (2) A reusable protocol for comparing cheap diagnostics with downstream performance under controlled degradation. (3) Practical guidance on when a good diagnostic score is a safe reason to skip real-data evaluation. A null result (the diagnostics react in step with mIoU) is also informative, since it supports using them as a shortcut.

\subsection{Provisional schedule}
\begin{table}[h]
\centering\small
\renewcommand{\arraystretch}{1.15}
\begin{tabularx}{\linewidth}{@{}l>{\raggedright\arraybackslash}X@{}}
\toprule
\textbf{Weeks} & \textbf{Activity}\\
\midrule
1--2 & Class mapping, data subset, degradation pipeline\\
3 & Diagnostics for all 11 datasets against real scans\\
4--6 & Training runs (11 datasets $\times$ 3 seeds) and evaluation\\
7 & Analysis of reaction points, per class and range\\
8 & Reporting and review\\
\bottomrule
\end{tabularx}
\end{table}
The weeks are relative and will be aligned with the confirmed course milestones.

\subsection{Collaboration}
Frans Laporte builds the degradation pipeline and computes the diagnostics. Florent Kegler sets up training and the per-class and per-range evaluation. Analysis, writing and review are shared, and each author checks the other's work.

\subsection{Risks and fallbacks}
\begin{itemize}
    \item \emph{Noisy or low synthetic-only baseline mIoU:} more seeds, the $\delta$ rule above, or optionally mixed synthetic and real training as an extension.
    \item \emph{Compute:} 33 training runs may be too many. Reduce the number of scans or epochs, use two seeds, or keep only one degradation form.
    \item \emph{Diagnostics depend on pretrained RangeNet++ and preprocessing:} use the released weights and settings unchanged.
    \item \emph{Class mapping mismatch:} restrict to classes that map cleanly and report the exclusions.
    \item \emph{Novelty not confirmed:} complete the documented literature search before the study starts.
\end{itemize}